\clearpage
\section*{Zusammenfassung}

Die Überwachung des Rutschhangs Schwanderbärgli oberhalb von Schwanden bei Brienz (BE) wird seit 1989 von der \gls{fhnw} 
durchgeführt. Die Folgemessung 2026 erfolgte im Feldkurs durch das Profil GeoSensorik \& Monitoring. Der vorliegende Bericht 
gibt einen Überblick über die durchgeführten GNSS-Messungen und die \gls{thb}. Mit einer Literaturrecherche wird der Einfluss 
des \gls{ztd} auf GNSS-Messungen beschrieben. Dabei wird gezeigt, dass die Berücksichtigung des \gls{ztd} in den GNSS-Messungen 
am Schwanderbärgli unumgänglich ist. Weiter wird dokumentiert, dass unterschiedliche Auswertestrategien zu teils signifikant 
anderen Resultaten der langstatischen GNSS-Messungen führen. Es ist somit unumgänglich, eine klare Vorgehensweise für die Durchführung 
und Auswertung der GNSS-Messungen am Schwanderbärgli zu definieren. Der Vorschlag für eine solche ist im Bericht beschrieben. 
Weiter zeigen die redundanten \acrshort{thb}-Messungen 2026 erstmals, dass das Anbringen eines Höhenmassstabes aus den Differenzen 
zwischen der \gls{thb} und den GNSS-Messungen kaum eine Steigerung der Genauigkeit bringt. Die Unsicherheiten in der \gls{thb} 
sind deutlich grösser als erwartet, zudem treten Widersprüche auf, deren Ursachen nicht abschliessend geklärt werden können. 
Folglich wird die \gls{thb} in der aktuellen Messkampagne nicht zur Korrektur der GNSS-Messungen eingesetzt, sondern fliesst 
als Beobachtung in das hybride Netz ein.

Die flächenhafte Deformationsmessung der Schwergewichtsmauer "Sperre" am Schwanderbärgli erfolgte mittels Laserscanning. Alle 
3D-Scans werden mittels M3C2-Algorithmus \parencite{lagueAccurate3DComparison2013} mit den früheren Epochen verglichen. 
Die dokumentierten Deformationen decken sich mit den Einzelpunktmessungen und betragen bis zu 15~cm im westlichen Bereich. 
Zu den Resultaten werden auch die Limitationen der verwendeten Methodik diskutiert. 

Unabhängig vom Schwanderbärgli wird ein Geomonitoring-Projekt am Campus Muttenz umgesetzt. Mit diesem wird der grundlegende 
Aufbau eines automatisierten Monitorings demonstriert und die Funktion der einzelnen Komponenten erklärt. Die aufgezeichneten 
Daten der Festpunkte über 24~h zeigen einen starken Drift in der Lage. Dieser wird von einer Teilkreisverdrehung verursacht 
und lässt sich in der Zeitreihenanalyse herausrechnen. Es wird gezeigt, dass die erhaltenen Daten durch eine Trendkorrektur 
ausreichend bereinigt werden können, um den Genauigkeitsanforderungen eines Monitorings zu genügen. 



